\section{Discussion and Conclusion}
\label{sec:discussion_conclusion}

This work frames sparse-sensor source apportionment as an identifiable-resolution
problem.  Known or proxy source inventories restrict the source space and make
attribution meaningful, but they do not guarantee that source groups are
separable from sparse observations.  The separability question is governed by
how inventory groups appear at the sensor network after wind-conditioned
transport, finite lag, background correction, and noise.

The central object is the projected lagged source-response matrix
\begin{equation}
\widetilde H=P_Q^\perp H_{1:T}^{\mathrm{lag}}.
\end{equation}
Exact source-level identifiability requires \(\operatorname{rank}(\widetilde H)=K\),
while robust attribution depends on the singular values, effective rank,
visibility, background absorption, pairwise coherence, and perturbation
sensitivity of \(\widetilde H\).  A model can fit sensor trajectories while
still producing unstable or non-unique source estimates, so prediction accuracy
alone is not evidence of source-level identifiability.

The practical implication is that source apportionment should report the finest
source grouping that the sensing system supports.  If two source groups have
highly coherent projected fingerprints, the data may support their combined
contribution but not their separate activities.  Reporting them separately can
create false precision.  The proposed IASA protocol therefore estimates source
activity, audits the projected response matrix, flags weak or ambiguous groups,
and merges source groups when the fine-grained inventory resolution is not
identifiable.

This perspective also changes how sensing systems and preprocessing choices
should be designed.  Adding sensors is useful when it improves visibility,
conditioning, or coherence of source fingerprints, not merely when it increases
spatial coverage.  Wind diversity can expose source groups from multiple
transport angles, while a single dominant wind regime can leave sources
confounded.  Background correction can remove broad temporal confounding, but a
basis that is too flexible can absorb source-driven variation and reduce
identifiability.  Wind interpolation and transport-model errors act as additional
observation error whose impact is amplified by ill-conditioning.

The framework has several limitations.  The response operator is an
approximation to atmospheric transport and can be extended with richer
meteorology, vertical mixing, deposition, chemistry, and boundary inflow.  The
theory is linear in source activities conditional on fixed wind, transport,
inventory, and background components; fully joint estimation introduces
additional nonlinearities.  Finally, the framework assumes that relevant
candidate source inventories are available.  Missing source categories may be
absorbed into available inventories or background terms, which should be treated
as a model-misspecification risk rather than as identifiable attribution.

Despite these limitations, the identifiability view provides a concrete basis
for more reliable policy-facing apportionment.  It separates fitting the observed
sensor signal from determining which source claims are supported by the sensing
system.  By using rank, singular values, visibility, background absorption,
coherence, and response perturbation sensitivity, source apportionment can report
not only estimated contributions, but also the resolution at which those
contributions can be trusted.

\section*{Acknowledgements}
Ankit Bhardwaj and Lakshminarayanan Subramanian were supported by the NSF Grant (Award Number 2335773) titled ``EAGER: Scalable Climate Modeling using Message-Passing Recurrent Neural Networks''. Lakshminarayanan Subramanian was also funded in part by the NSF Grant (award number OAC-2004572) titled ``A Data-informed Framework for the Representation of Sub-grid Scale Gravity Waves to Improve Climate Prediction''.
